% =========================================================
% Proof: Every Nonempty Finite Poset Has a Maximal Element
% =========================================================

\newpage
\phantomsection
\label{prf:finite-poset-maximal}

\begin{remark*}[Return]
\hyperref[prop:finite-poset-maximal]{Return to Proposition~\ref*{prop:finite-poset-maximal}.}
\end{remark*}

\begin{theorem*}[Every nonempty finite poset has a maximal element]
Let $(P, \leq)$ be a finite partially ordered set and let $Q \subseteq P$
be nonempty.  Then $\operatorname{Max}(Q) \neq \varnothing$.
\end{theorem*}

\begin{proof}[Proof]
Since $Q$ is nonempty, pick any $a_0 \in Q$.  If $a_0$ is maximal in $Q$,
we are done.  Otherwise there exists $a_1 \in Q$ with $a_0 < a_1$.
Continuing, if $a_1$ is not maximal there exists $a_2 \in Q$ with
$a_1 < a_2$.  The sequence $a_0 < a_1 < a_2 < \cdots$ is strictly
increasing in $Q$.  Since $Q \subseteq P$ and $P$ is finite, this sequence
cannot continue indefinitely.  Therefore it must terminate at some
$a_n \in Q$ from which no further strict increase is possible, and $a_n$
is maximal in $Q$.

For the second part, given any $x \in P$, apply the above argument to
$Q = \{y \in P : x \leq y\}$, which is nonempty (it contains $x$).  A
maximal element of $Q$ is an element of $\operatorname{Max}(P)$ that is
above $x$.
\end{proof}

\begin{remark*}[Proof structure]
An ascending chain argument.  The chain $a_0 < a_1 < \cdots$ must
terminate because $P$ is finite.  Finiteness is essential: the set of
all finite subsets of $\mathbb{N}$, ordered by inclusion, is infinite
and has no maximal element.
\end{remark*}

\begin{remark*}[Dependencies]
\hyperref[def:maximal-element]{Maximal element},
\hyperref[def:lattice-order-partial-order]{partial order (strict form)}.
\end{remark*}
